% 2 Objectives (all). Source: project plan section 2; deck slides 8-9.
\section{Objectives}\label{sec:objectives}

\paragraph{Broad objective} A diagnosis of \emph{whether}, \emph{where} and \emph{at what scale} AI data centres put
pressure on planetary boundaries, as a basis for responses matched to each scale of impact. The scale matters because
it decides who has to respond: a town council, a grid or river-basin authority, or an international agreement.

\paragraph{Specific objectives} Each objective covers one part of the Earth system and has one owner in the group.
\begin{description}[leftmargin=0pt,labelsep=0.4em,itemsep=1.5pt]
\item[O1 Energy and carbon] (atmosphere; Aditya). Map the regional distribution of data-centre electricity and grid
carbon intensity; compare growth rates across the three phases; explain why the cloud era stayed nearly flat while the
AI era does not.
\item[O2 Materials, water and waste] (hydrosphere and lithosphere; Aryan). Quantify the water-stress exposure of AI
capacity and of the people near it; estimate water use and identify water sources; trace the geography of critical
minerals and e-waste.
\item[O3 Economy, politics and society] (anthroposphere; Prerna). Measure how unevenly capacity is distributed
relative to population; document local opposition and supply-chain concentration.
\item[O4 Synthesis] (all). Place the findings on the planetary boundaries \citep{sakschewski2025} and the
safe-and-just Earth system boundaries \citep{rockstrom2023}, classified by scale of impact: local, regional or
global.
\end{description}
